% Source: https://artofproblemsolving.com/community/c3h1192155_perimeter
% Stable identity: aops-topic-1192155-q01
The diagonals bisect each other at the centre $O$, so the distance between a pair of opposite sides is twice the distance from $O$ to either of them. The two heights are therefore $4$ cm and $6$ cm.

Area $=$ base $\times$ height, so the sides are $\dfrac{24}{4}=6$ cm (the pair $4$ cm apart) and $\dfrac{24}{6}=4$ cm (the pair $6$ cm apart). The perimeter is $2(6+4)=\boxed{20}$ cm.

(In fact the parallelogram must be a rectangle: the sides of length $4$ join two lines that are $4$ apart, so they are perpendicular to them.)

\begin{asy}
usepackage("amsmath");
size(4.5cm);
defaultpen(fontsize(10pt));
pair A=(0,0),B=(6,0),C=(6,4),D=(0,4); pair O=(A+C)/2;
draw(A--B--C--D--cycle,linewidth(0.9)); draw(A--C,gray); draw(B--D,gray); dot(O);
draw(O--(3,0),dashed+blue); label("$2$",(3,1),E); draw(O--(0,2),dashed+red); label("$3$",(1.5,2),N);
label("$A$",A,SW);label("$B$",B,SE);label("$C$",C,NE);label("$D$",D,NW); label("$6$",(4.5,0),S); label("$4$",(6,2),E);
label("\mbox{}",(min(currentpicture,user=true)+max(currentpicture,user=true))/2,black); // keeps the text after the figure black when LaTeX typesets the labels inline
\end{asy}
